\documentclass[12pt, letterpaper]{article}

% ── Paquetes ──────────────────────────────────────────────────────────────
\usepackage[utf8]{inputenc}
\usepackage[T1]{fontenc}
\usepackage[spanish]{babel}
\usepackage{geometry}
\usepackage{amsmath, amssymb, amsthm}
\usepackage{booktabs}
\usepackage{array}
\usepackage{microtype}
\usepackage{setspace}
\usepackage{parskip}
\usepackage{titlesec}
\usepackage{fancyhdr}
\usepackage{enumitem}
\usepackage{bm}
\usepackage{mathtools}
\usepackage{tcolorbox}
\usepackage{pgfplots}
\pgfplotsset{compat=1.18}
\usepackage{tikz}
\usepackage{hyperref}

\tcbuselibrary{skins, breakable}

% ── Geometría ─────────────────────────────────────────────────────────────
\geometry{top=2.8cm, bottom=3.0cm, left=3.2cm, right=3.2cm}
\setstretch{1.20}

% ── Encabezado ────────────────────────────────────────────────────────────
\pagestyle{fancy}
\fancyhf{}
\fancyhead[C]{\small\textsc{Econometría \textperiodcentered{} Gujarati --- Ejercicio 10.5}}
\fancyfoot[C]{\small\thepage}
\renewcommand{\headrulewidth}{0.4pt}
\renewcommand{\footrulewidth}{0pt}

% ── Titulos de seccion ─────────────────────────────────────────────────────
\titleformat{\section}
  {\large\bfseries\scshape}{\thesection.}{0.6em}{}
  [\vspace{-4pt}\rule{\linewidth}{0.4pt}]
\titleformat{\subsection}{\normalsize\bfseries}{\thesubsection.}{0.5em}{}
\titleformat{\subsubsection}{\normalsize\itshape}{}{0em}{}

% ── Entornos ──────────────────────────────────────────────────────────────
\newtheorem{prop}{Proposicion}
\newtheorem{lema}{Lema}
\theoremstyle{remark}
\newtheorem{obs}{Observacion}

\newtcolorbox{clave}{
  colback=white, colframe=black,
  boxrule=0.5pt, arc=3pt,
  left=8pt, right=8pt, top=6pt, bottom=6pt,
  breakable
}

% ── Comandos ──────────────────────────────────────────────────────────────
\newcommand{\E}{\mathbb{E}}
\newcommand{\Var}{\operatorname{Var}}
\newcommand{\Cov}{\operatorname{Cov}}
\newcommand{\Corr}{\operatorname{Corr}}

% ─────────────────────────────────────────────────────────────────────────
\begin{document}

% ── Portada ───────────────────────────────────────────────────────────────
\begin{center}
  {\LARGE\bfseries Modelos de Rezago Distribuido}\\[6pt]
  {\LARGE\bfseries y Multicolinealidad en Series de Tiempo}\\[12pt]
  {\large\itshape Diagnostico, consecuencias y metodos de correccion}\\[14pt]
  \rule{0.55\linewidth}{0.6pt}\\[8pt]
  {\normalsize Ejercicio~10.5 --- \textit{Econometria}, Gujarati \& Porter, 5.a ed.}\\[4pt]
  {\small\textsc{Derivaciones completas: Koyck, Almon y primeras diferencias}}
\end{center}

\vspace{10pt}

\noindent\textbf{El modelo.}\quad
El modelo de rezago distribuido finito de orden 4 postula que el consumo
$Y_t$ en el periodo $t$ depende del ingreso corriente y de los cuatro periodos
anteriores:

\begin{equation}
  Y_t = \beta_1 + \beta_2 X_t + \beta_3 X_{t-1} + \beta_4 X_{t-2}
        + \beta_5 X_{t-3} + \beta_6 X_{t-4} + u_t
  \label{eq:modelo}
\end{equation}

\noindent donde $Y$ es el gasto de consumo, $X$ es el ingreso disponible y $t$
indexa el tiempo (trimestral). La interpretacion economica es que los hogares
ajustan su consumo de forma gradual ante cambios en el ingreso: el efecto total
de una variacion en $X$ se distribuye a lo largo de cinco periodos.

\medskip
\noindent\textbf{Parametros de interes:}
\begin{itemize}[leftmargin=1.6em, itemsep=3pt]
  \item $\beta_2$: efecto de impacto (impact multiplier) --- cambio inmediato en
    $Y$ ante un cambio unitario en $X_t$.
  \item $\beta_2 + \beta_3 + \cdots + \beta_6$: multiplicador de largo plazo
    (long-run multiplier) --- efecto acumulado total de un cambio permanente.
  \item $\beta_3, \ldots, \beta_6$: efectos rezagados en $t-1$ hasta $t-4$.
\end{itemize}

% ══════════════════════════════════════════════════════════════════════════
\section{Inciso (a): Multicolinealidad esperada}

\subsection{Respuesta}

\textbf{Si}. En modelos de rezago distribuido con datos de series de tiempo
economicas, la multicolinealidad severa es practicamente inevitable. El argumento
formal se desarrolla a continuacion.

\subsection{Autocorrelacion de series economicas como fuente de colinealidad}

Una serie es \textbf{estacionaria en covarianza} si:
\begin{align}
  \E[X_t] &= \mu \quad \text{(media constante)}\\
  \Var(X_t) &= \sigma^2_X \quad \text{(varianza constante)}\\
  \Cov(X_t, X_{t-s}) &= \gamma_s \quad \text{(covarianza que solo depende del rezago } s\text{)}
\end{align}

La \textbf{funcion de autocorrelacion} (FAC) en el rezago $s$ es:
\begin{equation}
  \rho_s = \frac{\Cov(X_t, X_{t-s})}{\Var(X_t)} = \frac{\gamma_s}{\sigma^2_X}
  \label{eq:autocorr}
\end{equation}

Para series de ingreso trimestral es comun que $\rho_1 > 0.9$, lo que significa
que $X_t, X_{t-1}, X_{t-2}, X_{t-3}, X_{t-4}$ estan altamente correlacionadas
entre si.

\subsection{Estructura de la matriz de correlaciones}

La matriz de correlaciones entre los regresores de~( \ref{eq:modelo} ) tiene la
estructura simetrica de Toeplitz:

\begin{equation}
  \mathbf{R} =
  \begin{pmatrix}
    1      & \rho_1  & \rho_2  & \rho_3  & \rho_4  \\
    \rho_1 & 1       & \rho_1  & \rho_2  & \rho_3  \\
    \rho_2 & \rho_1  & 1       & \rho_1  & \rho_2  \\
    \rho_3 & \rho_2  & \rho_1  & 1       & \rho_1  \\
    \rho_4 & \rho_3  & \rho_2  & \rho_1  & 1
  \end{pmatrix}
  \label{eq:matriz_R}
\end{equation}

\noindent La entrada $(j,k)$ es $\rho_{|j-k|}$. Si $\rho_1 \to 1$, todas las
entradas fuera de la diagonal se aproximan a~1, la matriz se aproxima a
$\mathbf{1}\mathbf{1}'$ (singular) y $\det(\mathbf{R}) \to 0$.

\subsection{Cuantificacion: VIF en funcion de la autocorrelacion}

En el caso simplificado de dos regresores $X_t$ y $X_{t-1}$ con correlacion
$\rho_1$, el $R^2$ auxiliar es $\rho_1^2$ y por tanto:

\begin{equation}
  \text{VIF}(X_t) = \frac{1}{1 - \rho_1^2}
  \label{eq:VIF_rho}
\end{equation}

\begin{table}[ht]
\centering
\caption{VIF como funcion de la autocorrelacion $\rho_1$ (dos rezagos)}
\label{tab:VIF}
\small
\begin{tabular}{cccc}
\toprule
$\rho_1$ & $\rho_1^2$ & $1 - \rho_1^2$ & VIF \\
\midrule
0.50 & 0.250 & 0.750 & 1.33 \\
0.70 & 0.490 & 0.510 & 1.96 \\
0.80 & 0.640 & 0.360 & 2.78 \\
0.90 & 0.810 & 0.190 & 5.26 \\
0.95 & 0.903 & 0.098 & 10.2 \\
0.99 & 0.980 & 0.020 & 50.0 \\
\bottomrule
\end{tabular}
\end{table}

\noindent Con cinco regresores, los VIF son aun mayores porque cada $X_{t-s}$
esta correlacionado con \emph{todos} los demas rezagos simultaneamente.

\subsection{Consecuencias practicas}

Bajo multicolinealidad severa en~(\ref{eq:modelo}):

\begin{enumerate}[leftmargin=1.8em, itemsep=4pt]
  \item \textbf{Estimadores inestables}: pequenas variaciones en la muestra
    producen cambios grandes en $\hat\beta_3, \ldots, \hat\beta_6$.

  \item \textbf{Errores estandar inflados}: la varianza de cada $\hat\beta_j$
    escala con $\text{VIF}_j$, reduciendo los estadisticos $t$:
    \begin{equation}
      \Var(\hat\beta_j) = \frac{\sigma^2}{S_{jj}(1 - R_j^2)} = \sigma^2 \cdot
      \frac{\text{VIF}_j}{S_{jj}}
      \label{eq:var_vif}
    \end{equation}

  \item \textbf{Signos incorrectos}: es posible obtener $\hat\beta_4 < 0$
    (el ingreso de hace dos periodos \emph{reduce} el consumo), lo cual carece
    de sentido economico. La colinealidad impide separar el efecto individual
    de cada rezago.

  \item \textbf{Alto $R^2$ con estadisticos $t$ no significativos}: el ajuste
    global es bueno (el ingreso en sus rezagos explica bien el consumo) pero
    ningun coeficiente individual supera el umbral de significancia.
\end{enumerate}

\begin{clave}
  Con $\rho_1 > 0.9$ y cinco rezagos de la misma variable, los VIF superan
  10 para todos los coeficientes. Las estimaciones individuales de
  $\hat\beta_3, \ldots, \hat\beta_6$ son practicamente inutiles para inferencia.
\end{clave}

% ══════════════════════════════════════════════════════════════════════════
\section{Inciso (b): Metodos de correccion}

\subsection{Transformacion de Koyck}

\subsubsection{Supuesto central}

Koyck propone que los coeficientes de los rezagos declinan geometricamente:
\begin{equation}
  \beta_{s+2} = \beta_2 \lambda^s, \qquad s = 0, 1, 2, \ldots
  \qquad 0 < \lambda < 1
  \label{eq:koyck_supuesto}
\end{equation}

El modelo de rezago infinito bajo este supuesto es:
\begin{equation}
  Y_t = \beta_1 + \beta_2 \sum_{s=0}^{\infty} \lambda^s X_{t-s} + u_t
  \label{eq:koyck_infinito}
\end{equation}

\subsubsection{Derivacion algebraica}

Rezagamos~(\ref{eq:koyck_infinito}) un periodo y multiplicamos por $\lambda$:
\begin{equation}
  \lambda Y_{t-1} = \lambda\beta_1 + \beta_2\lambda \sum_{s=0}^{\infty}\lambda^s X_{t-1-s} + \lambda u_{t-1}
  \label{eq:koyck_rezagado}
\end{equation}

Restando~(\ref{eq:koyck_rezagado}) de~(\ref{eq:koyck_infinito}):
\begin{align}
  Y_t - \lambda Y_{t-1}
  &= \beta_1(1-\lambda)
     + \beta_2\left(\sum_{s=0}^{\infty}\lambda^s X_{t-s}
     - \lambda\sum_{s=0}^{\infty}\lambda^s X_{t-1-s}\right)
     + u_t - \lambda u_{t-1}
  \label{eq:koyck_resta1}
\end{align}

Expandimos la diferencia de sumas:
\begin{align}
  \sum_{s=0}^{\infty}\lambda^s X_{t-s}
  - \lambda\sum_{s=0}^{\infty}\lambda^s X_{t-1-s}
  &= X_t + \sum_{s=1}^{\infty}\lambda^s X_{t-s}
   - \sum_{s=1}^{\infty}\lambda^s X_{t-s} = X_t
  \label{eq:cancelacion}
\end{align}

Los terminos en $X_{t-s}$ para $s \geq 1$ se cancelan exactamente. Sustituyendo:
\begin{equation}
  \boxed{
    Y_t = \underbrace{\beta_1(1-\lambda)}_{\alpha_0}
          + \underbrace{\beta_2}_{\alpha_1} X_t
          + \underbrace{\lambda}_{\alpha_2} Y_{t-1}
          + v_t
  }
  \label{eq:koyck_final}
\end{equation}
donde $v_t = u_t - \lambda u_{t-1}$ es un error de media movil de orden~1
(MA(1)). El modelo tiene ahora \textbf{solo tres parametros} en lugar de seis.

\subsubsection{Recuperacion de parametros}

Dados $\hat\alpha_0, \hat\alpha_1, \hat\alpha_2$:
\begin{align}
  \hat\beta_2 &= \hat\alpha_1\\
  \hat\lambda &= \hat\alpha_2\\
  \hat\beta_1 &= \frac{\hat\alpha_0}{1-\hat\alpha_2}
\end{align}

El \textbf{multiplicador de largo plazo} es:
\begin{equation}
  M_{\infty} = \frac{\hat\beta_2}{1-\hat\lambda} = \frac{\hat\alpha_1}{1-\hat\alpha_2}
  \label{eq:mult_LP}
\end{equation}

\subsubsection{Problema econometrico}

La transformacion introduce $Y_{t-1}$ como regresor y $v_t = u_t - \lambda u_{t-1}$
como error. Si $u_t$ es autocorrelacionado, $Y_{t-1}$ y $v_t$ estaran
correlacionados, haciendo a MCO \textbf{inconsistente}. La solucion correcta
es usar variables instrumentales (VI) o minimos cuadrados en dos etapas (MC2E).

\subsection{Rezagos distribuidos polinomiales: metodo de Almon}

\subsubsection{Supuesto central}

En lugar del decaimiento geometrico de Koyck, Almon supone que los coeficientes
$\beta_2, \ldots, \beta_6$ siguen un \textbf{polinomio de grado $p$} en el
indice del rezago $s$:

\begin{equation}
  \beta_{s+2} = a_0 + a_1 s + a_2 s^2 + \cdots + a_p s^p,
  \qquad s = 0, 1, 2, 3, 4
  \label{eq:almon_supuesto}
\end{equation}

\subsubsection{Derivacion de las variables transformadas}

Sustituyendo~(\ref{eq:almon_supuesto}) en~(\ref{eq:modelo}):
\begin{align}
  Y_t
  &= \beta_1 + \sum_{s=0}^{4}\beta_{s+2} X_{t-s} + u_t \notag\\
  &= \beta_1 + \sum_{s=0}^{4}\left(\sum_{j=0}^{p} a_j s^j\right) X_{t-s} + u_t \notag\\
  &= \beta_1 + \sum_{j=0}^{p} a_j \underbrace{\left(\sum_{s=0}^{4} s^j X_{t-s}\right)}_{Z_{jt}} + u_t
  \label{eq:almon_derivacion}
\end{align}

Para $p = 2$ las variables transformadas son:
\begin{align}
  Z_{0t} &= X_t + X_{t-1} + X_{t-2} + X_{t-3} + X_{t-4}
  \label{eq:Z0}\\
  Z_{1t} &= 0\cdot X_t + 1\cdot X_{t-1} + 2\cdot X_{t-2} + 3\cdot X_{t-3} + 4\cdot X_{t-4}
  \label{eq:Z1}\\
  Z_{2t} &= 0\cdot X_t + 1\cdot X_{t-1} + 4\cdot X_{t-2} + 9\cdot X_{t-3} + 16\cdot X_{t-4}
  \label{eq:Z2}
\end{align}

El modelo estimable es:
\begin{equation}
  \boxed{Y_t = \beta_1 + a_0 Z_{0t} + a_1 Z_{1t} + a_2 Z_{2t} + u_t}
  \label{eq:almon_final}
\end{equation}

\noindent con \textbf{tres parametros} ($a_0, a_1, a_2$) en lugar de seis.

\subsubsection{Recuperacion de los coeficientes originales}

Dados $\hat a_0, \hat a_1, \hat a_2$:
\begin{align}
  \hat\beta_2 &= \hat a_0 \label{eq:rec0}\\
  \hat\beta_3 &= \hat a_0 + \hat a_1 + \hat a_2 \label{eq:rec1}\\
  \hat\beta_4 &= \hat a_0 + 2\hat a_1 + 4\hat a_2 \label{eq:rec2}\\
  \hat\beta_5 &= \hat a_0 + 3\hat a_1 + 9\hat a_2 \label{eq:rec3}\\
  \hat\beta_6 &= \hat a_0 + 4\hat a_1 + 16\hat a_2 \label{eq:rec4}
\end{align}

En forma matricial, $\hat{\boldsymbol{\beta}}_{2:6} = \mathbf{P}\hat{\mathbf{a}}$ con:
\begin{equation}
  \mathbf{P} =
  \begin{pmatrix}
    1 & 0 & 0 \\ 1 & 1 & 1 \\ 1 & 2 & 4 \\ 1 & 3 & 9 \\ 1 & 4 & 16
  \end{pmatrix}
  \label{eq:almon_P}
\end{equation}

\subsubsection{Comparacion con Koyck}

\begin{itemize}[leftmargin=1.6em, itemsep=3pt]
  \item \textbf{Flexibilidad}: el polinomio acomoda perfiles no monotonas
    (efecto maximo en $t-2$, decreciente despues).
  \item \textbf{Sin endogeneidad}: no se introduce $Y_{t-1}$.
  \item \textbf{Eleccion del grado}: $p$ debe ser menor que el numero de
    rezagos. Para 4 rezagos, $p=4$ recupera exactamente MCO sin restriccion
    (sin ganancia). Se elige $p$ por AIC.
\end{itemize}

\subsection{Primeras diferencias}

Si la colinealidad proviene de una tendencia estocastica comun (series
integradas I(1)), la transformacion en primeras diferencias puede eliminarla.
Definiendo $\Delta X_t = X_t - X_{t-1}$:

\begin{equation}
  \Delta Y_t = \beta_2 \Delta X_t + \beta_3 \Delta X_{t-1}
               + \beta_4 \Delta X_{t-2} + \beta_5 \Delta X_{t-3}
               + \beta_6 \Delta X_{t-4} + \Delta u_t
  \label{eq:diferencias}
\end{equation}

\subsubsection{Por que las diferencias tienen menor correlacion}

Sea $X_t = \phi X_{t-1} + \varepsilon_t$ con $|\phi| < 1$ (AR(1)). La
correlacion entre diferencias consecutivas es:

\begin{equation}
  \rho_{\Delta,1}
    = \frac{\Cov(\Delta X_t, \Delta X_{t-1})}{\Var(\Delta X_t)}
    = \frac{-\phi}{1 + \phi^2}
  \label{eq:corr_diff}
\end{equation}

Para $\phi = 0.9$: $\rho_{\Delta,1} = -0.9/1.81 \approx -0.497$. Las
diferencias tienen correlacion negativa de magnitud mucho menor que la
correlacion positiva $\rho_1 = 0.9$ en niveles.

\subsubsection{Limitacion}

Si las series estan cointegradas, el modelo en diferencias omite el
mecanismo de correccion de error:
\begin{equation}
  \Delta Y_t = \alpha(Y_{t-1} - \hat\mu X_{t-1})
               + \sum_{s=0}^{4}\beta_{s+2}\Delta X_{t-s} + u_t
  \label{eq:mce}
\end{equation}
perdiendo informacion de largo plazo.

\subsection{Simplificacion secuencial (ad hoc)}

Se elimina secuencialmente el regresor con menor estadistico $t$ hasta que
todos sean significativos. Este metodo adolece de \textbf{sesgo por pretesting}:
si $\beta_6 \neq 0$ pero su $t$ es pequeno (por colinealidad, no por nulidad),
eliminarlo introduce sesgo por variable omitida:

\begin{equation}
  \E[\hat\beta_j^{\text{red}}]
    = \beta_j + \sum_{k \notin \text{modelo}} \beta_k \cdot \gamma_{kj}
  \label{eq:sesgo_omision}
\end{equation}

Con colinealidad alta, $\gamma_{kj}$ puede ser grande, amplificando el sesgo.
\textbf{No se recomienda para inferencia formal.}

% ══════════════════════════════════════════════════════════════════════════
\section{Comparacion de metodos}

\begin{table}[ht]
\centering
\caption{Comparacion de metodos para multicolinealidad en rezago distribuido}
\label{tab:comparacion}
\small
\begin{tabular}{p{2.4cm}p{3.8cm}p{3.8cm}p{2.6cm}}
\toprule
\textbf{Metodo} & \textbf{Ventajas} & \textbf{Desventajas} & \textbf{Cuando usar} \\
\midrule
Koyck &
Parsimonia (3 parametros); multiplicador LP directo &
Endogeneidad de $Y_{t-1}$; decaimiento monotono forzado &
Ajuste dinamico rapido \\
\midrule
Almon (PDL) &
Flexible; sin endogeneidad; grado controlable &
Eleccion de $p$ arbitraria &
Efecto diferido con pico interior \\
\midrule
Primeras diferencias &
Elimina tendencia &
Pierde relacion de largo plazo &
Series I(1) sin cointergracion \\
\midrule
Ad hoc &
Facil de implementar &
Sesgo por pretesting &
Solo exploracion inicial \\
\bottomrule
\end{tabular}
\end{table}

% ══════════════════════════════════════════════════════════════════════════
\section{Sintesis}

\begin{enumerate}[leftmargin=1.8em, itemsep=6pt]
  \item La multicolinealidad es \textbf{esperada y severa} en~(\ref{eq:modelo})
    porque el ingreso trimestral tiene autocorrelacion tipicamente mayor a
    $0.9$, produciendo VIF superiores a 10.

  \item El metodo preferido es \textbf{Almon} cuando el perfil de rezagos es
    no monotono, y \textbf{Koyck} (con correccion por endogeneidad via VI)
    cuando se espera decaimiento geometrico.

  \item Las \textbf{primeras diferencias} son apropiadas solo si las series
    son I(1) y no cointegradas.

  \item La \textbf{simplificacion ad hoc} puede usarse como diagnostico
    inicial pero nunca para inferencia formal.
\end{enumerate}

% ══════════════════════════════════════════════════════════════════════════
\appendix

\section{Derivacion completa de la transformacion de Koyck}

Partimos de:
\begin{equation}
  Y_t = \mu + \beta\sum_{s=0}^{\infty}\lambda^s X_{t-s} + u_t
  \label{eq:kA}
\end{equation}

Rezagamos y multiplicamos por $\lambda$:
\begin{equation}
  \lambda Y_{t-1} = \lambda\mu + \beta\lambda\sum_{s=0}^{\infty}\lambda^s X_{t-1-s} + \lambda u_{t-1}
  \label{eq:kB}
\end{equation}

Restando (\ref{eq:kB}) de (\ref{eq:kA}) y expandiendo:
\begin{align}
  Y_t - \lambda Y_{t-1}
  &= \mu(1-\lambda) + \beta\left(X_t
     + \sum_{s=1}^{\infty}\lambda^s X_{t-s}
     - \sum_{s=1}^{\infty}\lambda^s X_{t-s}\right)
     + u_t - \lambda u_{t-1} \notag\\
  &= \mu(1-\lambda) + \beta X_t + v_t
\end{align}

Despejando $Y_t$:
\begin{equation}
  Y_t = \mu(1-\lambda) + \beta X_t + \lambda Y_{t-1} + v_t \qquad \square
\end{equation}

\section{Construccion matricial de las variables de Almon}

Para $p=2$ y cuatro rezagos, la relacion entre los $Z_{jt}$ y los rezagos de
$X$ es:

\begin{equation}
  \begin{pmatrix} Z_{0t} \\ Z_{1t} \\ Z_{2t} \end{pmatrix}
  =
  \begin{pmatrix}
    1 & 1 & 1 & 1 & 1 \\
    0 & 1 & 2 & 3 & 4 \\
    0 & 1 & 4 & 9 & 16
  \end{pmatrix}
  \begin{pmatrix} X_t \\ X_{t-1} \\ X_{t-2} \\ X_{t-3} \\ X_{t-4} \end{pmatrix}
  = \mathbf{P}'\,\mathbf{x}_t
\end{equation}

La estimacion por MCO de $\mathbf{a}$ en el modelo transformado equivale a
la estimacion restringida de $\boldsymbol{\beta}$ bajo la restriccion de que
los coeficientes caigan sobre el polinomio~(\ref{eq:almon_supuesto}).

\section{Perfil tipico de coeficientes}

\begin{figure}[ht]
\centering
\begin{tikzpicture}
\begin{axis}[
  width=9cm, height=5cm,
  xlabel={Rezago $s$},
  ylabel={Coeficiente $\hat\beta_{s+2}$},
  xtick={0,1,2,3,4},
  xticklabels={$t$, $t{-}1$, $t{-}2$, $t{-}3$, $t{-}4$},
  ymin=-0.05, ymax=0.55,
  grid=major, grid style={dashed, black!20},
  axis lines=left,
  legend pos=north east,
  legend style={font=\small},
  label style={font=\small},
  tick label style={font=\small},
]
\addplot[mark=square*, black, thick, dashed]
  coordinates {(0,0.50)(1,0.30)(2,0.18)(3,0.11)(4,0.06)};
\addlegendentry{Koyck ($\lambda=0.6$)}
\addplot[mark=*, black, thick]
  coordinates {(0,0.10)(1,0.28)(2,0.38)(3,0.30)(4,0.14)};
\addlegendentry{Almon cuadratico}
\end{axis}
\end{tikzpicture}
\caption{Koyck: decaimiento monotono. Almon: permite pico en $t-2$.}
\end{figure}

\vspace{14pt}
\noindent\rule{\linewidth}{0.4pt}

\noindent{\small\textit{Nota.} Los valores numericos de las figuras son
ilustrativos. El modelo de Koyck con $Y_{t-1}$ requiere estimacion por VI o
MC2E para ser consistente en presencia de autocorrelacion serial en $u_t$.}

\end{document}
